% 这个顾虑很成熟，而且比"堆满高级引用"要明智得多。本科thesis的评分逻辑是**"你是否真的理解你用的每一个概念、并能讲清它和你工作的关系"**——一个你能完全讲透的概念，比十个你只是引用了名字的概念更能证明你的能力。真正露怯的不是"用了简单概念"，而是"用了一个高级概念但明显没吃透"。所以你的取舍方向完全对。

% 我帮你把这些文献分成三档：**必须讲透的、点到为止的、可以砍掉的**。

% ## 第一档：必须讲透，因为它们是你工作的地基

% **Set partitioning —— 这个不能砍，而且它恰恰是最"配得上你能力"的一个**

% 你担心的方向反了。set partitioning不是一个会让你"显得超纲"的高级概念，它是一个**让你的形式化站得住脚的基础概念**，而且它直接回答了你评委一定会问的那个问题：**"你为什么用这个分组方法（全局DP），而不是更直觉的聚类或贪心？"**

% 你的论证链是这样的，每一步你都能亲口讲清楚：
% 1. Connections的规则是"16个词分成4个互斥的组，每个词恰好用一次"
% 2. 这个结构在数学上有个名字，叫**集合划分问题**——把一个集合切成若干互斥子集、每个元素恰属一个、优化某目标
% 3. 因为它是集合划分问题，所以"每组恰好4个、不重不漏"是一个**硬约束**，而普通聚类（KMeans）不保证这个约束
% 4. 所以我用全局DP做精确求解，而不是聚类或贪心——**因为问题的结构要求全局约束满足**
% 5. 实验①证明了这个选择是对的（全局0.48 > 贪心0.34）

% **看到没——set partitioning不是装点门面，它是你"为什么这么设计"的理由。** 你用它，是因为你需要一个词来精确描述你面对的问题结构，从而解释你的方法选择。这是本科thesis里最该有的东西：**方法选择有理有据，而不是"我试了这个正好work"**。评委看到这个会觉得"这个学生理解自己在做什么"，而不是"这个学生在炫技"。

% 而且它一点都不难讲透——定义两句话，和Connections的对应关系一目了然，你完全hold得住。**留着，认真讲，这是加分项不是风险项。**

% **Wordle的信息论对照 —— 留着，因为你是真的用它解释了你的负结果**

% 你说"从Wordle获得灵感很好"——比灵感更强，你实际上是用Wordle**解释了你为什么失败**。你的bit分析（Wordle反馈7.9 bits/次 vs 你的1.58 bits/次）不是掉书袋，它是你explore负结果的**因果解释**：Wordle的信息论探索能work，是因为反馈带宽大；你的不work，是因为带宽只有它的1/5。

% 这个对照**配得上你的能力**，因为它展示的不是"我知道Wordle"，而是"我能用一个类比问题来诊断我自己方法的局限"。这是本科生能做到的、很漂亮的分析。留着，但讲的时候聚焦在那个bit数字的对比上，不需要展开Wordle求解器的完整技术细节——你只需要它的反馈带宽这一个数字。

% ## 第二档：点到为止，一句话带过，不展开

% **Knuth Mastermind** —— 用一句话定位就够，不要展开算法细节。

% 写法建议：*"这种'提交猜测—根据反馈缩小可能性—再猜测'的序贯结构，与经典的Mastermind解密游戏同构；Knuth(1977)最早为其给出了系统的求解策略。"* ——一句话，给你的交互式框架一个学术出处，然后就走。**不要**去详述minimax算法、不要证明5步可解——那些是Mastermind论文的内容，不是你的，展开了反而像在充数。你只需要它证明"我的问题不是我瞎想的，它属于一个有历史的问题族"。

% **Dr.Fill** —— 同样一句话，作为"谜题当CSP解"的先例。

% 写法建议：*"将谜题形式化为约束满足问题并非新思路，Ginsberg(2011)的Dr.Fill填字求解器即将填字游戏建模为加权CSP。"* ——点到,证明你的形式化思路有先例、不是孤立的，然后走。你没实现填字求解器，所以不需要深入它的技术。

% ## 第三档：可以砍掉，或降到脚注

% **Combinatorial group testing** —— 我建议**砍掉，或最多一个脚注**。

% 理由：你**没有真正实现**一个group testing算法（你的endgame分支被贝叶斯滤波自然吸收了）。引用一个你没实际使用的理论框架，正是你担心的那种"显得超纲"的风险——评委可能会问"你既然提了group testing，那你的算法和标准group testing算法怎么比？"而你没有答案，因为你没实现它。**引用一个你答不上后续问题的概念，是真正的风险。** 除非你想在理论讨论里专门写那个endgame的组合分析（那需要额外的篇幅和把握），否则砍掉最干净。如果实在舍不得，放一个脚注："剩余8词的endgame可视为一个组合群测试问题，本文中它被贝叶斯滤波自然涵盖，故不单独处理。"——一句话交代，不展开。

% ## 一个总的原则，帮你之后自己判断

% **判断一个概念该不该留、该讲多深，问自己一个问题：如果评委指着它问"这个和你的工作具体什么关系？为什么用它？"，我能不能不看稿子、用自己的话讲两分钟？**

% - 能（set partitioning、Wordle带宽对照）→ 讲透，这是你的强项
% - 只能讲一句它是什么、但讲不出深入关系（Knuth、Dr.Fill）→ 一句话点到，作为出处
% - 讲不清它和你工作的实际关系（group testing，因为你没实现）→ 砍掉或脚注

% **这个原则能保护你**：你论文里出现的每个概念，都是你能扛得住追问的。这样整篇论文的"知识水位"就和你的真实理解匹配了——不是靠砍掉难概念来匹配，而是靠"只留下你真懂的"来匹配。而你会发现，真正撑起你论文的那几个（集合划分、全局约束满足、bit带宽、先验主导），你其实全都懂，因为它们是你一路做实验做出来的，不是读来的。

% 这才是本科thesis最好的状态：**你论文里的每一句理论，背后都有你自己的一个实验或一次推导撑着。** 那样它就不可能"和你的能力不匹配"——因为它就是你能力的产物。


\chapter{Related Work}
\label{ch:related_work}

This chapter provides an overview of prior work relevant to this thesis, including work that directly solves Connections using LLMs and work that inspired the methods used in this study. It first introduces prior work on solving Connections with AI. It then presents related work on the methods used, which serves as the knowledge base for this study.

\section{Solving Connections with AI}

Current work on this topic mainly uses LLMs as the main solution for solving Connections. These studies compare how well different models solve Connections. Some use Connections as a benchmark to test LLM capability, while others explore different prompting techniques and methods. \citeauthor{toddMissedConnectionsLateral2024} use three methods, namely a sentence-embedding baseline, a comparison between two LLMs (GPT-3.5-Turbo-1106 and GPT-4-1106-Preview), and \ac{CoT} prompting. Their experiments allow five incorrect and five invalid guesses per puzzle, which departs from the original game's rules. GPT-4 Turbo with \ac{CoT} achieves the highest success rate at 58.11\%. The best-performing sentence-embedding baseline (MPNet) reaches 11.6\%. LLMs are often stumped by categories that involve non-semantic word properties, abstract features, or context-dependent usage. They also report a finding relevant to this thesis. Puzzle order matters, since presenting puzzles in their original, grouped order achieves substantially higher accuracy than presenting them in randomized order~\cite{toddMissedConnectionsLateral2024}.

A similar work from \citeauthor{samadarshiConnectingDotsEvaluating2024} benchmarked several LLMs against both novice and expert human players and creates a taxonomy of knowledge types required to solve connections. They used in total 438 Connections puzzles and their best performing models Claude Sonnet 3.5 achieved 18\% success rate of fully solving the puzzles. They categorizes NYT Connections categories along two main dimensions, word form (covering morphology, orthography, phonology, and multiword expressions) and word meaning (covering semantic relations such as synonymy, polysemy, and hypernymy, along with associative and encyclopedic relations), with a further category for groups that combine both form-based and meaning-based cues. Results show that models performed comparably well on semantic relation categories but significantly worse on the other types.

More recent work from \citeauthor{loredolopezNYTConnectionsDeceptivelySimple2025} compare three game configurations: single attempt, four trys without feedback (one away prompt) and four trys with feedback. Human solvers outperform LLM under all three configurations. The best performing model Claude 3.5 under full hints reaches 40.25\%, while the human full-hints score reaches 60.67\%. They also experiment on different prompting techniques. \ac{CoT} and self-consistency prompting yield diminishing or even negative returns as puzzle difficulty increases. Simple input-output prompting sometimes outperforms these more elaborate reasoning strategies. Their herustic (non-LLM) baseline, using Multilingual-E5-Large-Instruct embedding and beam search, reaches 13.25\% accuracy.

\citeauthor{zhangLearningSemanticComplexities} compare in-context prompting of a large model (DBRX) against fine-tuning and distilling a smaller model (Mistral~7B) on solved puzzles. They found that fine-tuning a smaller model (Mistral 7B) directly on solved puzzles fails to capture the task itself, which underperformed even their clustering baseline. Their strongest result comes from Verbal Reinforcement Learning, in which a larger supervisor model critiques the fine-tuned model's output in natural language across three rounds, roughly doubling the average per-puzzle match rate relative to single-pass prompting \cite{zhangLearningSemanticComplexities}.

% set partitioning

% wordle: infromation theory

